% Notes of the Montreal seminar, 1970 (cote 129, typed part): archive cote 129, site page 50 (PDF page 51).
% typescript. Draft compared: none.
% AI-transcribed by Claude Opus 5.5 (claude-opus-5-5, adaptive thinking, effort high), Message Batches API, batch mtl129_full (msgbatch_01AGnDKHwy6ATCM7ehUa6263), 2026-11-13; not yet checked by eye.
\scanpage{129/50}
\mspage{\hand{II.22}}

adjoint à gauche de $D^*$ , il résultera du fait que $D^*$ est une équivalen-%
ce de catégories que c'est un foncteur quasi-inverse de $D^*$ .

6.1. Soit $M$ un $\hand{\mathcal{D}}$-module de longueur finie sur $W$ et où $V$ est
nilpotent sur $M$ . On lui associe le foncteur en groupes $\mathbb{E}(M)$ sur
(sch/k) défini par
\[ \mathbb{E}(M)(S) = \mathrm{Hom}_{\hand{\mathcal{D}}}(M, \underset{\hand{\to}}{\mathbb{W}}(S)) \]
pour tout k-schéma $S$ .

Il faut démontrer que $\mathbb{E}(M)$ est représenté par un p-groupe
fini unipotent et que le foncteur $\mathbb{E}$ ainsi défini est adjoint à gauche
de $D^*$ .

6.2. $\mathbb{E}(M)$ \uline{est représentable par un schéma affine}

On choisit un entier $N$ suffisamment grand pour que $V^{N+1}$
soit nul sur $M$ et on considère les polynômes $P_N$ et $S_N$ qui définis-%
sent respectivement la dernière coordonnée de la multiplication et de
l'addition dans $\mathbb{W}_{N+1}$ (ce sont des polynômes sur $\mathbb{Z}$ à $2N+2$ variables).

On définit alors la k-algèbre $A$ comme quotient de
$k[T_x, x \in M]$ par les relations suivantes :
\[ \text{i)}\quad T_{F(x)} = T_x^p \]
\[ \text{ii)}\quad T_{x+y} = S_N(T_{V^N(x)}, \ldots, T_x \,;\, T_{V^N(y)}, \ldots, T_y) \]
\margin{$\times$}
\[ \text{iii)}\quad T_{\hand{\lambda}x} = P_N(\hand{\lambda}_1^{p^{1-N}}, \hand{\lambda}_2^{p^{1-N}}, \ldots, \hand{\lambda}_{N+1}^{p^{1-N}} \,;\, T_{V^N(\hand{x})}, \ldots, T_{\hand{x}}) \ . \]

Pour prouver que $\mathrm{Spec}(A)$ représente le foncteur $\mathbb{E}(M)$ ,
on remarque que pour tout $x \in M$ et tout $u \in \mathrm{Hom}_{\hand{\mathcal{D}}}(M, \underset{\hand{\to}}{\mathbb{W}}(S))$ , on a

\mspage{\hand{50}}
